
\subsubsection*{5.2.3}
\phantomsection
\label{sga3:VIIB:5.2.3}

Suppose now that \(B\) has finite height at most \(r\).
Let
\[
  \pi:k\xrightarrow{\ \sim\ }k,
  \qquad
  x\longmapsto x^p,
\]
be the Frobenius automorphism. The linear map
\[
  B\widehat\otimes_{\pi}k\longrightarrow B,
  \qquad
  b\widehat\otimes_{\pi}x
  \longmapsto
  b^px=(bx^{1/p})^p,
\]
has closed image equal to the closed subalgebra
\[
  B^p=\{\,b^p\mid b\in B\,\}
\]
of \(B\). Put
\[
  J=B^p\cap I=B^p\cap I_A.
\]

\footnote{\label{sga3:VIIB:nde:158}The editors added what follows.}
Let \(G_1\) be the kernel of the Frobenius morphism
\[
  \operatorname{Fr}:G\longrightarrow G^{(p)},
\]
and let \(HG_1\) be the formal subgroup of \(G\) that is the
inverse image of the formal subgroup \(H^{(p)}\subset G^{(p)}\).
Then \(HG_1\) is defined by the closed ideal generated by the
\(p\)-th powers of the elements of \(AI\); this ideal is \(AJ\).
Moreover, the formation of \(G/H\) commutes with base change,
because \(G/H\) represents the quotient sheaf for the flat
topology \textup{(cf.\ \SGARefLink{sga3:VIIB:2.4}{2.4})}. Hence
\[
  (G/H)^{(p)}=G^{(p)}/H^{(p)},
\]
and we have the following commutative diagrams:

\begin{figure}[H]
  \centering
  \begin{tikzcd}[row sep=2.5em,column sep=2.8em]
G \arrow[r,"\operatorname{Fr}"] \arrow[d]
  & G^{(p)} \arrow[d] \\
G/H \arrow[r,"\operatorname{Fr}"']
  & G^{(p)}/H^{(p)}
\end{tikzcd}
\qquad
\begin{tikzcd}[row sep=2.5em,column sep=2.8em]
A
  & A\widehat\otimes_\pi k
      \arrow[l,"{a\widehat\otimes1\mapsto a^p}"'] \\
B \arrow[u]
  & B\widehat\otimes_\pi k
      \arrow[l,"{b\widehat\otimes1\mapsto b^p}"']
      \arrow[u]
\end{tikzcd}

\end{figure}

It follows that \(B^p\) is the affine algebra of the quotient
\(G/HG_1\).\footnote{\label{sga3:VIIB:nde:159}As indicated in the
original, this also follows from \SGARefLink{sga3:VIIB:5.1:proposition}{Proposition~5.1}. The editors
preferred to give the preceding argument, which does not use the
implication \textup{\SGARefLink{sga3:vii:paragraph:5-1:equation:i}{(i)}} \(\Rightarrow\) \textup{\SGARefLink{sga3:VIIB:5.1:proposition:item:ii}{(ii)}} of that
proposition.}
Temporarily denote by \(C\) the affine algebra of the quotient
\(HG_1/H\). Since the formation of \(G/H\) commutes with base
change, the square

\begin{figure}[H]
  \centering
  \begin{tikzcd}[row sep=2.8em,column sep=3.5em]
HG_1 \arrow[r] \arrow[d]
  & G \arrow[d] \\
HG_1/H \arrow[r]
  & G/H
\end{tikzcd}

\end{figure}

\noindent
is Cartesian. Consequently,
\[
  A\widehat\otimes_BC
  \simeq
  A/AJ
  =
  A\widehat\otimes_B(B/BJ).
\]
Since \(A\) is topologically free over \(B\)
\textup{(by 2.4, because \(A\) and \(B\) are local)}, the natural
morphism
\[
  B/BJ\longrightarrow C
\]
is an isomorphism. Thus \(B/BJ\) is the affine algebra of
\(HG_1/H\), and of course
\[
  B/BJ=B_1
\]
has height at most \(1\), since
\[
  J=\bigl((x^p)\bigr)_{x\in\mathfrak r(B)}.
\]

Let
\[
  B'
  =
  k[[\omega]]
  \big/
  \bigl((x^{p^r})\bigr)_{x\in\omega_r},
\]
let \(\phi:B'\to B\) be the morphism introduced in \SGARefLink{sga3:VIIB:5.2.1:section-proposition}{5.2.1}, let
\[
  B'^p=\{\,x^p\mid x\in B'\,\},
\]
and let \(J'\) be the augmentation ideal of \(B'^p\). We then
have the commutative diagram

\begin{figure}[H]
  \centering
  \begin{tikzcd}[row sep=2.8em,column sep=3.7em]
B' \arrow[r,"\psi"] \arrow[d]
  & B \arrow[d] \\
B'_1=B'/\overline{B'J'}
  \arrow[r,"\phi_1"]
  & B_1=B/\overline{BJ}
\end{tikzcd}

\end{figure}

\noindent
and, by \SGARefLink{sga3:VIIB:5.2.2}{5.2.2}, \(\phi_1\) is an isomorphism.

On the other hand, by 2.4, \(A\) is topologically flat over
\[
  B=\mathscr A(G/H)
  \qquad\text{and over}\qquad
  B^p=\mathscr A(G/HG_1).
\]
It follows from \SGARefLink{sga3:VIIB:1.3.3}{1.3.3} that \(B\) is topologically flat over
\(B^p\). Moreover, \SGARefLink{sga3:VIIB:5.2.4}{5.2.4} below shows that the morphism
\[
  B'^p\longrightarrow B^p
\]
induced by \(\phi\) is an isomorphism. We may therefore apply
\SGARefLink{sga3:VIIB:0.3.4}{0.3.4} to the pseudocompact ring \(B'^p=B^p\) and the
pseudocompact \(B^p\)-modules \(M=B'\), \(N=B\). By what
precedes,
\[
  \phi_1
  =
  \phi\widehat\otimes_{B^p}k
\]
is an isomorphism, and it follows that \(\phi\) itself is an
isomorphism. This proves \SGARefLink{sga3:VIIB:5.2:theorem}{5.2} when \(B\) has finite height, subject
only to \SGARefLink{sga3:VIIB:5.2.4}{5.2.4} below.
